\chapter{How a huge file scrolls}\label{ch:huge}
\begin{goals}
\textbf{In this chapter} you learn how lamina reads a file it never loads: the index it builds, how a
screenful is read, how filters, sorts and statistics pass through the file, and what each of them costs.
It is the chapter to read before you open a file larger than your memory.
\end{goals}

\Cref{sec:windowed} gave the idea in a paragraph. This chapter gives the mechanism, because knowing it
tells you what will be fast, what will take a pass through the file, and where the limits are.

\section{The block index}\label{sec:block-index}

When you open a delimited file, lamina reads it once from start to end. It doesn't keep what it reads.
It counts records (a record is a line, unless a quoted field carries a line break into the next), and
every 4{,}096 records it notes the byte offset where the next one starts. The list of those offsets is the
\term{block index}. With it, the file is a sequence of blocks of 4{,}096 records each, and lamina knows
where each block starts.

The index is small: one number per block, eight bytes. A file of 50~million rows has about 12{,}200
blocks, an index of about 100~KB. That, and a few blocks of the file, is all lamina holds for the file
itself.

\begin{center}
\begin{tikzpicture}[x=1mm, y=1mm, font=\sffamily\scriptsize, text=ink]
  \foreach \i in {0,...,9} { \fill[tint] (\i*12, 0) rectangle ++(11.2, 7); }
  \fill[ambersoft] (36, 0) rectangle ++(11.2, 7); \fill[ambersoft] (48, 0) rectangle ++(11.2, 7);
  \foreach \i in {0,...,9} { \draw[amber, line width=0.6pt] (\i*12, 7) -- ++(0, 3); }
  \node[anchor=west] at (0, 13) {block index: where every 4{,}096th record starts};
  \node[anchor=west, text=muted] at (0, -4) {the file, in blocks of 4{,}096 records};
  \draw[{Stealth[length=2mm]}-, sel, line width=0.6pt] (48, -1) -- (62, -9) node[anchor=west, text=sel] {the blocks on screen: read, decoded, kept};
  \node[anchor=west, text=muted] at (121, 3.5) {…};
\end{tikzpicture}
\end{center}

\subsection{Building it}

The pass that builds the index runs in a background worker, so the window stays responsive and you can
watch the status line count \code{indexing… 37\%}. It reads the file in chunks straight from the disk, so
a file larger than memory is indexed without being loaded. Where a worker can't be started (lamina on
\file{file://}, \cref{sec:single-file}), the same pass runs on the page itself, and the status line
counts it the same way.

The detection of the file's format (\cref{sec:detection}) uses only its first 64~KB and happens before
the pass, so a misread file can be corrected before you wait for a full index.

\subsection{The index cache}\label{sec:index-cache}

lamina keeps each file's index on your machine, keyed by the file's name, size and modification time
(\cref{sec:local-state}). Open the same file again and the index comes from the cache: the table is
there without a pass. Change the file, and its size or modification time changes with it, so it is
indexed afresh. The cache holds offsets, not data.

\section{Reading a screenful}\label{sec:screenful}

To show rows 1{,}204{,}000 to 1{,}204{,}040, lamina finds the block they are in (row number divided by
4{,}096), reads that block's bytes from the file (from its offset to the next one's), splits it into
records and decodes the ones on screen. The last 16 blocks read are kept, so scrolling within them, or
back to them, reads nothing. Jumping anywhere is the same single read, which is why the row box answers
at once at any depth.

\section{Scans}\label{sec:scans}

A filter, a sort, a statistic or an export has to look at every row, so it is a \term{scan}: one pass
through the file, block by block, keeping only what the operation needs. While it runs, the status line
counts its progress, and \key{Esc} cancels it, leaving the view as it was.

\begin{description}
\item[A filter] keeps, for each matching row, where it is in the file: its byte offset and length, and
its row number. Twenty-four bytes per match, however wide the row. The filtered view then reads each
row on screen directly from its offset.
\item[A sort] keeps, for each row, the same location plus its sort keys, then orders them. Sorting a
filtered view sorts only its matches. This is why a large file is best filtered before it is sorted.
\item[Statistics] keep running sums (count, mean, the sums of squared and cubed deviations) and, for
the exact quantiles, the column's values as numbers (\cref{sec:one-column}).
\item[Group by and grade–tonnage] keep one set of running sums per group.
\item[An export] keeps nothing: each row is written as it is read (\cref{ch:export}).
\end{description}

Scans compose. A sort runs over the current filter's matches; statistics, group by and grade–tonnage
run over them too. Every instrument respects the filter.

\section{The record cursor}\label{sec:cursor}

Underneath, every format reaches the scans the same way, through a \term{record cursor}: something that
can walk the records in order and re-read one by its location. For a CSV the location is a byte offset
and length. For a \file{.dm} it is the record number. For a Parquet file it is the row's number in the
file. The filter, the sort, the statistics and the export are written once, against the cursor, and
work on every format unchanged. That is why a \file{.dm} filters and sorts exactly as a CSV does.

\subsection{Strided reads of a .dm}\label{sec:strided}

A \file{.dm} record has the same length everywhere and each field sits at the same place inside it. To
read one field from every record, lamina can step through the file in strides of the record length and
decode just that field, without decoding the rest. A filter uses this: it reads only the columns its
expression names. On a wide model and a filter over one or two columns, that is most of the work saved.

\subsection{Column reads of a Parquet file}

Parquet stores each column of a row group separately, so a scan reads only the columns it needs, and a
filter also skips whole row groups its footer statistics rule out (\cref{sec:pruning}).

\begin{note}
Both of these are about the filter's own columns. When calculated columns exist, a filter on a
\file{.dm} or Parquet file reads every column, because a calculated column may use any of them
(\cref{ch:calc}).
\end{note}

\section{What it costs}\label{sec:costs}

\begin{center}\small
\begin{tabularx}{\linewidth}{@{}lL@{}}
\toprule
\sffamily\bfseries you do & \sffamily\bfseries it costs \\
\midrule
open a file the first time & one pass, to build the index \\
open it again & nothing: the index is cached \\
scroll, jump to a row, resize & reading one block, or none \\
filter & one pass; 24 bytes per match \\
sort & one pass over the rows (or the matches); their locations and keys in memory \\
statistics of one column & one pass; the column's values, for the quantiles \\
column summary (Σ stats) & two passes, every column at once \\
group by, grade–tonnage & one pass each (grade–tonnage's curves, a second) \\
export & one pass; nothing kept \\
the distributions over the headers & a sample of 8{,}192 rows, read when the file opens \\
\bottomrule
\end{tabularx}
\end{center}

A pass costs what reading the file costs: the disk's speed, plus decompression for a compressed source,
plus decoding the fields the scan needs. A filter over a CSV decodes every row (the text has to be split
to find any field); a filter over a \file{.dm} or Parquet decodes only its columns.

\section{The limits}\label{sec:huge-limits}

Holding only locations has a price: the number of locations. A filter keeps at most 16{,}777{,}216
matches and a sort at most 5{,}242{,}880 rows. Past the first, the filter stops with
\code{too many matches — refine the filter} and the view stays as it was; past the second, the sort
stops with \code{too many rows to sort — filter first} and the view is shown unsorted. The exact quantiles of a column are computed over at most 5{,}242{,}880 values; past
that, the statistics window shows the rest without them. \Cref{app:limits} lists every limit.

\begin{tip}
On a very large file, filter first. A filter is cheap in memory and makes every later sort, statistic
and export smaller. \code{Z between 400 and 500} on a model of tens of millions of blocks gives a bench
you can sort and summarise freely.
\end{tip}

\section{Resident sources}\label{sec:resident}

Some sources aren't windowed: they are decompressed or decoded into memory whole. A workbook, an archive
of up to 32~MB, a Zstandard, xz or bzip2 file, the sample block model. They behave exactly as a windowed
file does, only the whole thing is in memory, so they are limited by it, which is why lamina refuses the
compressed kinds above 32~MB (\cref{sec:archives}). A gzip file or zip entry above 32~MB is read through
the tape instead (\cref{sec:tape}).
